\documentclass[tikz,border=4pt]{standalone}
\usepackage{ctex}
\usepackage{amsmath}
\usepackage{pgfplots}
\pgfplotsset{compat=1.18}

% p10-02b: 一般式、顶点式、交点式三种形式对比示意
% y=x^2-2x-3（一般式） vs y=(x-1)^2-4（顶点式）
\begin{document}
\begin{tikzpicture}
\begin{axis}[
  width=8cm, height=8cm,
  axis lines=center,
  xlabel={$x$}, ylabel={$y$},
  every axis x label/.style={at={(axis cs:4.5,0)}, anchor=west},
  every axis y label/.style={at={(axis cs:0,3.0)}, anchor=south},
  xmin=-2.0, xmax=4.5,
  ymin=-4.8, ymax=3.0,
  xtick={-2,-1,1,2,3,4},
  ytick={-4,-3,-2,-1,1,2},
  tick style={color=black},
  ticklabel style={font=\small},
  clip=false,
]
  % y = x^2 - 2x - 3（蓝色）
  \addplot[blue, thick, domain=-1.5:3.5, samples=100] {x^2 - 2*x - 3};
  \node[right, blue, font=\small] at (axis cs:3.5, 1.5) {$y=x^2-2x-3$};

  % 顶点式标注 y=(x-1)^2-4（注释在顶点处）
  \addplot[only marks, mark=*, mark size=3pt, red] coordinates {(1,-4)};
  \node[below, red] at (axis cs:1,-4) {$(1,-4)$};
  \node[below right, red, font=\small] at (axis cs:1.2,-3.5) {顶点式：$y=(x-1)^2-4$};

  % 零点 (-1,0) 和 (3,0)（交点式的因式）
  \addplot[only marks, mark=*, mark size=3pt] coordinates {(-1,0) (3,0)};
  \node[above] at (axis cs:-1,0) {\small$-1$};
  \node[above] at (axis cs:3,0) {\small$3$};
  \node[above, font=\small] at (axis cs:1, 0.5) {交点式：$y=(x+1)(x-3)$};

  % 对称轴
  \draw[gray, dashed] (axis cs:1,-4.8) -- (axis cs:1,3.0);
  \node[above, gray, font=\small] at (axis cs:1,3.0) {$x=1$};
\end{axis}
\end{tikzpicture}
\end{document}
